\documentclass{article}
\usepackage[margin=1.5cm]{geometry}
\usepackage{amsmath}
\usepackage{enumitem}

\setlist{nosep,leftmargin=*}

\begin{document}
\twocolumn

\title{Chapter 10 Warm-Up: Dictionaries and a Phone Book}
\author{Prof. Jordan C. Hanson}
\date{}

\maketitle

\section{Dictionary Basics}

\begin{enumerate}

\item A Python dictionary stores information as \textbf{key--value pairs}. Consider

\begin{verbatim}
phone_book = {
 "562-555-0142": "Maria Garcia",
 "562-555-0188": "Daniel Kim",
 "562-555-0125": "Elena Garcia",
 "562-555-0171": "Maria Lopez",
 "562-555-0194": "Aisha Patel"
}
\end{verbatim}

Each phone number is a \textbf{key}.  The corresponding name is its \textbf{value}.  What is printed by

\begin{verbatim}
print(phone_book["562-555-0188"])
\end{verbatim}

\vspace{0.1cm}

Why is it useful for each dictionary key to be unique?

\end{enumerate}


\section{Creating and Adding Entries}

\begin{enumerate}

\item An empty dictionary can be created using

\begin{verbatim}
phone_book = {}
\end{verbatim}

A new key--value pair can then be added:

\begin{verbatim}
phone_book["562-555-0142"] = "Maria Garcia"
\end{verbatim}

Add another entry:

\begin{verbatim}
phone_book["562-555-0188"] = "Daniel Kim"
\end{verbatim}

Print the entire dictionary:

\begin{verbatim}
print(phone_book)
\end{verbatim}

What happens if you assign a new value to a key that already exists?

\end{enumerate}


\section{Dictionary Operations}

\begin{enumerate}

\item The method \texttt{phone\_book.keys()} returns the dictionary keys. The method \texttt{phone\_book.values()} returns the values. The method \texttt{phone\_book.items()} provides the key and value together.  For our phone book \texttt{(phone, name)} represent one key--value pair.

\end{enumerate}


\section{Searching by Name}

\begin{enumerate}

\item We can combine dictionaries with the string-searching
techniques from Chapter 8.  Begin with

\begin{verbatim}
search = input("Enter first/last name, or both: ")
\end{verbatim}

Then examine each entry:

\begin{verbatim}
for phone, name in phone_book.items():
    if search.lower() in name.lower():
        print(name, phone)
\end{verbatim}

Predict the output when the user enters

\begin{verbatim}
Garcia
\end{verbatim}

How many phone numbers should be printed?

\end{enumerate}


\section{Put the Search in a Function}

\begin{enumerate}

\item Write a function that returns all matching phone numbers:

\begin{verbatim}
def find_numbers(phone_book,search):
    matches = []
    ...
    return matches
\end{verbatim}

Test the function:

\begin{verbatim}
print(find_numbers(phone_book,"Maria"))
\end{verbatim}

Which phone numbers should be returned?

\end{enumerate}


\section{Warm-Up Exercise}

Using the original five-entry dictionary, predict the result of
each search before running the program:

\begin{verbatim}
find_numbers(phone_book,"Garcia")
find_numbers(phone_book,"Maria")
find_numbers(phone_book,"Daniel Kim")
find_numbers(phone_book,"Johnson")
\end{verbatim}

\noindent
Answer the following:
\begin{itemize}
\item Which search returns two phone numbers?
\item Which search uses both a first and last name?
\item What does the function return when there are no matches?
\item Why is a list used for \texttt{matches} instead of returning the first phone number immediately?
\end{itemize}

\end{document}